% The electron distance on Lewis-labeled graphs.
% Companion note to paper/lwg/main.tex. Standalone; compiles on its own.
\documentclass[11pt]{article}

\usepackage{amsmath,amssymb,amsfonts}
\usepackage{amsthm}
\usepackage{booktabs}
\usepackage{longtable}
\usepackage{array}
\usepackage{geometry}
\geometry{margin=1.05in}
\usepackage[version=4]{mhchem}
\usepackage{hyperref}
\hypersetup{colorlinks=true,linkcolor=blue,citecolor=blue,urlcolor=blue}

\numberwithin{equation}{section}

\newcommand{\Loc}{\operatorname{Loc}}
\newcommand{\gout}{\operatorname{out}}
\newcommand{\gin}{\operatorname{in}}
\newcommand{\Ie}{I_{\mathrm e}}
\newcommand{\ValG}{\operatorname{Val}}
\newcommand{\lp}{\operatorname{lp}}
\newcommand{\del}{\operatorname{d}}
\newcommand{\ch}{\operatorname{ch}}
\newcommand{\dsc}{\operatorname{dsc}}
\newcommand{\Occ}{\mathfrak{C}}
\newcommand{\LLG}{\mathfrak{L}}
\newcommand{\code}[1]{\texttt{#1}}
\DeclareMathOperator{\supp}{supp}

\theoremstyle{plain}
\newtheorem{theorem}{Theorem}[section]
\newtheorem{proposition}[theorem]{Proposition}
\newtheorem{lemma}[theorem]{Lemma}
\newtheorem{corollary}[theorem]{Corollary}
\theoremstyle{definition}
\newtheorem{definition}[theorem]{Definition}
\theoremstyle{remark}
\newtheorem{observation}{Observation}
\newtheorem{remark}[theorem]{Remark}
\newtheorem{example}[theorem]{Example}

\title{The Electron Distance on Lewis-Labeled Graphs:\\
Integrality, Parity, and Composition}
\author{Companion note to\\
\emph{Lewis-labeled graphs: curly arrows and fishhooks
as executable electron transfers}}
\date{\today}

\begin{document}
\maketitle

\begin{abstract}
Chemical distance is classically measured by counting the bonds that a
reaction forms and breaks.  On a Lewis-labeled graph (LLG) the natural
refinement also counts nonbonding electrons, weighting a lone pair like a
bond-order unit and an unpaired electron by one half.  We show that the
resulting quantity $\del$ is the weighted $\ell^1$ metric on the vector of
locus occupancies, and that it admits a complete arithmetic theory.  It is
integer-valued on every multiset of electron moves, admissible or not.  It
equals the number of electrons displaced, hence $2c+f$ for $c$ curly arrows
and $f$ fishhooks in the absence of cancellation.  On integral move multisets
its parity is governed by the radical signature: $\del\equiv r^+\equiv
r^-\pmod2$, where $r^\pm$ count the unpaired electrons created and destroyed.
Two consequences delimit the polar and the radical case exactly.  Groups of
two-electron moves are integral automatically and satisfy $\del=2|M|$, so
polar chemistry is even unconditionally and at every level of description.
Integral radical groups are even for the signatures $(0,0)$, $(2,0)$, $(0,2)$
and $(2,2)$---closed shell, initiation, termination---and odd for $(1,1)$,
chain propagation; $\beta$-scission has $\del=3$.  Parity is not additive
along a mechanism: the defect is the number of chain-carrier round trips, so
a radical chain may have odd step sum and even net distance.  The identity
$\sum_v\Delta r_v\equiv0$ recovers Eq.~(4.3) of Holzschuh \emph{et al.}~[1] as
the degenerate case.  All statements are verified against the complete set of
worked examples, grammar classes and reviewed corpus records of the companion
manuscript.
\end{abstract}

\tableofcontents

\section{Introduction}
\label{sec:intro}

Let a chemical reaction be given together with an atom-to-atom map.  The
oldest quantitative summary of such a reaction is its \emph{chemical
distance}, the number of bonds formed plus the number of bonds broken.  On the
attributed graphs of classical reaction informatics this is the only
invariant available, because those graphs do not represent nonbonding
electrons.  On a Lewis-labeled graph, where lone pairs and unpaired electrons
are stored fields on the same footing as $\sigma$- and $\pi$-bond orders, the
definition can be refined: one counts a change of one lone pair like a change
of one bond-order unit, and a change of one unpaired electron by one half.
The refinement is forced by the arithmetic rather than chosen, since these are
exactly the weights that make each unit worth two electrons.

The purpose of this note is to develop the resulting quantity as a
mathematical object and to determine, completely, when it is even.  The
question is not idle.  A curly arrow moves two electrons and a fishhook moves
one; if the distance counts electrons, then the parity of the distance records
whether the mechanism is built from pairs alone.  This is what one expects to
be the formal residue of the classical distinction between polar and radical
chemistry, and it is.  We show that groups of two-electron moves are even
unconditionally, that integral radical groups are even or odd according to a
signature that separates initiation and termination from propagation, and that
parity fails to be additive along a radical chain in a way that is itself
chemically meaningful.

Throughout, the notation of the companion manuscript is used without
comment, and the half-edge graph formalism of Holzschuh, Flamm, Merkle and
Stadler~[1] is used for comparison.  Their setting and ours describe the same
electronic bookkeeping at different levels of resolution, and several of the
statements below have recognizable antecedents there; these are identified
where they occur.

\section{Preliminaries}
\label{sec:prelim}

\subsection{Lewis-labeled graphs and electron loci}

Recall from the companion manuscript that a \emph{Lewis-labeled graph} is an
attributed undirected graph $G=(V,E,a,b)$ in which each atom carries
$a(v)=(el_v,\ell_v,r_v)$ and each edge carries $b(e)=(s_e,p_e)$, with
$o_e=s_e+p_e$ the total Kekul\'e bond order.  Its \emph{electron loci} are
\begin{equation}
  \Loc(G)=\{\lp(v),\rho_v: v\in V\}\cup
          \{\sigma_{uv},\pi_{uv}: \{u,v\}\in\tbinom{V}{2}\},
  \label{eq:loci}
\end{equation}
with occupancies $c_G(\lp(v))=\ell_v$, $c_G(\rho_v)=r_v$,
$c_G(\sigma_{uv})=s_{uv}$, $c_G(\pi_{uv})=p_{uv}$, extended by zero on
non-edges.  Each locus has a sort in $\{\lp,\rho,\sigma,\pi\}$ and an
\emph{electron quotum}
\begin{equation}
  w(\lp)=w(\sigma)=w(\pi)=2,\qquad w(\rho)=1 ,
  \label{eq:quota}
\end{equation}
the number of electrons one unit of occupancy represents.  An \emph{electron
move} is a triple $\mu=(x\xrightarrow{k}y)$ with $x\neq y$ loci and
$k\in\{1,2\}$; its occupancy delta is $\Delta_\mu=-\tfrac{k}{w(x)}\mathbf1_x
+\tfrac{k}{w(y)}\mathbf1_y$, and for a finite multiset $M$ of moves we write
$\Delta_M=\sum_{\mu\in M}\Delta_\mu$, together with the gross outflow and
inflow $\gout_M$, $\gin_M$, so that $\Delta_M=\gin_M-\gout_M$.  A move with $k=2$ is
a \emph{curly arrow}, one with $k=1$ a \emph{fishhook}.  We write
\begin{equation}
  c=\#\{\mu\in M: k_\mu=2\},\qquad f=\#\{\mu\in M: k_\mu=1\}
  \label{eq:arrow-counts}
\end{equation}
for their numbers.  The multiset $M$ is \emph{integral} if
$\Delta_M\in\mathbb Z^{\Loc(G)}$; the admissibility conditions of the
companion manuscript---locality, common-prestate availability, integrality and
$\Pi$-validity of the successor---are recalled only where they are used.
Fix a finite atom set $V$ and write
\begin{equation}
  \Occ(V)=\mathbb N_0^{\Loc(V)}
  \label{eq:occ-space}
\end{equation}
for the set of occupancy vectors on it.  This is the ambient space of all
statements below; it plays the role of the set $\mathfrak H[H,V]$ of half-edge
graphs on a fixed electron and atom set in~[1].

\subsection{Electron accounting and its half-edge antecedent}

It is convenient to record the occupancy delta in electron units.

\begin{definition}
\label{def:electron-delta}
The \emph{electron delta} of a move multiset $M$ is the function
$E_M\colon\Loc(G)\to\mathbb Z$,
\begin{equation}
  E_M(x)=w(x)\,\Delta_M(x) .
  \label{eq:electron-delta}
\end{equation}
\end{definition}

That $E_M$ takes values in $\mathbb Z$ and not merely in $\tfrac12\mathbb Z$ is
the content of Lemma~\ref{lem:E-integral} below; it is the first point at
which the quota~\eqref{eq:quota} does work.

\begin{lemma}
\label{lem:conservation}
For every finite multiset $M$ of electron moves,
\begin{equation}
  \sum_{x\in\Loc(G)}E_M(x)=0 .
  \label{eq:conservation}
\end{equation}
\end{lemma}

\begin{proof}
A move $x\xrightarrow{k}y$ contributes $w(x)\bigl(-k/w(x)\bigr)
+w(y)\bigl(+k/w(y)\bigr)=-k+k=0$, and~\eqref{eq:conservation} is the sum of
these contributions.  No hypothesis on $M$ beyond finiteness is used.
\end{proof}

Equation~\eqref{eq:conservation} is the conservation lemma of the companion
manuscript.  In the half-edge picture it corresponds to Eq.~(4.2) of~[1],
$\dsc_\alpha(\varnothing)+\sum_{v}\dsc_\alpha(v)=0$, where
$\dsc_\alpha(v)=|\psi'(\alpha(v))|-|\psi(v)|$ measures the change in the number
of half-edges associated with $v$.  The two identities express the same fact,
that the number of electrons cannot change, but they are indexed differently:
Eq.~(4.2) of~[1] distributes electrons over \emph{atoms}, whereas
\eqref{eq:conservation} distributes them over \emph{loci}, so that a bond
appears once rather than at each of its two ends.  The locus indexing is what
makes the parity results of Section~\ref{sec:parity} available.

\begin{observation}
\label{obs:heg-dict}
Under the correspondence of~[1, Eqs.~(2.2)--(2.3)] one has $L(v)=\ell_v$,
$R(v)=r_v$ and $B(v)=\sum_{e\sim v}o_e$, so that $\deg(v)=B(v)+2L(v)+R(v)$
becomes $\deg(v)=n_v+\sum_{e\sim v}o_e$ with $n_v=2\ell_v+r_v$, and the formal
charge $\ch(v)=-(\deg(v)-\mathrm{val}(a(v)))$ of~[1, Eq.~(2.4)] coincides with
the derived charge $q_v$ of the companion manuscript.  The dictionary is used
below only to locate antecedents; no result depends on it.
\end{observation}

\section{The electron distance}
\label{sec:distance}

\subsection{Definition}

\begin{definition}
\label{def:distance}
The \emph{electron distance} of a finite multiset $M$ of electron moves is
\begin{equation}
  \del(M)=\tfrac12\sum_{x\in\Loc(G)}\bigl|E_M(x)\bigr|
        =\tfrac12\sum_{x\in\Loc(G)}w(x)\bigl|\Delta_M(x)\bigr| .
  \label{eq:distance}
\end{equation}
Written out in the four label fields,
\begin{equation}
  \del(M)=\sum_{e}\Bigl(|\Delta s_e|+|\Delta p_e|\Bigr)
         +\sum_{v}|\Delta\ell_v|
         +\tfrac12\sum_{v}|\Delta r_v| ,
  \label{eq:distance-fields}
\end{equation}
so that a bond-order unit and a lone pair each count $1$ and an unpaired
electron counts $\tfrac12$.  The classical, electron-blind quantity is the
bond term alone, aggregated per edge,
\begin{equation}
  \del_{\mathrm b}(M)=\sum_e\bigl|\Delta s_e+\Delta p_e\bigr|=\sum_e|\Delta o_e| .
  \label{eq:distance-bond}
\end{equation}
\end{definition}

\begin{example}
\label{ex:ammonia-hcl}
For $\ce{NH3 + HCl -> NH4+ + Cl-}$, the two curly arrows
$\lp(N)\to\sigma_{NH}$ and $\sigma_{HCl}\to\lp(Cl)$ give $\Delta\ell_N=-1$,
$\Delta s_{NH}=+1$, $\Delta s_{HCl}=-1$ and $\Delta\ell_{Cl}=+1$.  The
electron-blind reading records one bond formed and one broken,
$\del_{\mathrm b}=2$, while $\del=1+1+1+1=4$.  The discrepancy is exactly the
pair of nonbonding changes that a bond-only account cannot see.
\end{example}

The two readings of the bond term in~\eqref{eq:distance-fields}
and~\eqref{eq:distance-bond} agree under the policy of the companion
manuscript.

\begin{observation}
\label{obs:sigma-pi}
If both endpoint states satisfy $\Pi_{\mathrm{exec}}$ then
$|\Delta s_e|+|\Delta p_e|=|\Delta s_e+\Delta p_e|$ for every edge $e$.
Indeed $s_e\in\{0,1\}$ forces $\Delta s_e\in\{-1,0,1\}$.  If $\Delta s_e=-1$
then $s_e$ ends at $0$, and $p_e>0\Rightarrow s_e=1$ forces the final
$p_e$ to vanish, whence $\Delta p_e\le0$.  If $\Delta s_e=+1$ then $s_e$
starts at $0$, and the same implication forces the initial $p_e$ to vanish,
whence $\Delta p_e\ge0$.  The two components never move in opposite
directions.
\end{observation}

\subsection{The metric property}

Definition~\ref{def:distance} depends on $M$ only through $\Delta_M$, and is
therefore a function on differences of occupancy vectors.

\begin{proposition}
\label{prop:metric}
The assignment
\begin{equation}
  \del(c,c')=\tfrac12\sum_{x\in\Loc(V)}w(x)\,\bigl|c'(x)-c(x)\bigr|
  \label{eq:metric}
\end{equation}
is a metric on $\Occ(V)$.  It is the $\ell^1$ metric for the weight system
$\tfrac12 w$, which assigns weight $1$ to every pair-valued locus and
$\tfrac12$ to every radical locus.  If $M$ transforms $c$ into $c'$, then
$\del(M)=\del(c,c')$.
\end{proposition}

\begin{proof}
The weights $\tfrac12 w(x)$ are strictly positive, so~\eqref{eq:metric} is a
positively weighted $\ell^1$ norm applied to $c'-c$; symmetry, positivity, the
identity of indiscernibles and the triangle inequality are inherited.  The
final assertion is Definition~\ref{def:distance} with $\Delta_M=c'-c$.
\end{proof}

\begin{remark}
\label{rem:vs-editing-distance}
Proposition~\ref{prop:metric} should be contrasted with the editing distance
of~[1, \S3.1], defined as the minimum number of editing operations---bond
breakings, bond formations and half-edge movements---converting one half-edge
graph into another, and shown there to be a metric on $\mathfrak H[H,V]$ by
invertibility of the individual operations (Lemma~3.2 of~[1]) together with
concatenation.  The two quantities are of different type.  The editing
distance is a minimum over realizing sequences and counts \emph{operations};
$\del$ is a closed-form function of the endpoints and counts \emph{electrons}.
The relation between them is Proposition~\ref{prop:arrow-count} below: $\del$
is a lower bound for the electron cost of any realizing sequence of pushes,
attained exactly when the sequence performs no cancellation.  Since [1,
Prop.~3.3] guarantees that a realizing sequence of electron pushes always
exists, $\del$ is the minimum electron cost of an explanation, not merely a
bound.
\end{remark}

\section{Integrality}
\label{sec:integrality}

\begin{lemma}
\label{lem:E-integral}
$E_M(x)\in\mathbb Z$ for every finite multiset $M$ and every locus $x$, and
$\sum_x|E_M(x)|$ is even.
\end{lemma}

\begin{proof}
A move contributes $\pm k$ to $E_M$ at each of its two loci, since
$w(x)\cdot k/w(x)=k$, and $k\in\{1,2\}$; thus $E_M$ is a finite sum of
integers.  By Lemma~\ref{lem:conservation} and $|a|\equiv a\pmod 2$ applied
termwise, $\sum_x|E_M(x)|\equiv\sum_x E_M(x)=0\pmod2$.
\end{proof}

\begin{theorem}
\label{thm:integer}
$\del(M)\in\mathbb N_0$ for every finite multiset $M$ of electron moves, and
$\del(M)=0$ if and only if $\Delta_M=0$.  No admissibility hypothesis is
required.
\end{theorem}

\begin{proof}
Immediate from Lemma~\ref{lem:E-integral} and~\eqref{eq:distance}.
\end{proof}

\begin{remark}
\label{rem:integrality-weaker}
Theorem~\ref{thm:integer} is strictly weaker than integrality of $M$ and must
not be confused with it.  A single fishhook $\sigma_{uv}\to\rho_u$ has
$\Delta(\sigma_{uv})=-\tfrac12\notin\mathbb Z$ and is rejected by the
integrality gate, yet $E(\sigma_{uv})=-1$, $E(\rho_u)=+1$ and $\del=1$.  The
electron distance is defined, and integral, on move multisets that describe no
transition at all.  Section~\ref{sec:diagnostics} isolates what it does
detect.
\end{remark}

\begin{proposition}
\label{prop:arrow-count}
Let $M$ be a finite multiset of electron moves with arrow counts~$c$ and~$f$
as in~\eqref{eq:arrow-counts}.  Then
\begin{equation}
  \sum_{\mu\in M}k_\mu-\del(M)
  =\sum_{x\in\Loc(G)}w(x)\min\bigl(\gin_M(x),\gout_M(x)\bigr)\;\ge\;0 .
  \label{eq:arrow-count-defect}
\end{equation}
In particular, if no locus is both a source and a target of moves in $M$, then
\begin{equation}
  \del(M)=\sum_{\mu\in M}k_\mu=2c+f .
  \label{eq:arrow-count}
\end{equation}
\end{proposition}

\begin{proof}
Summing $w(x)\gout_M(x)=\sum_{\mu\text{ from }x}k_\mu$ over $x$ gives
$\sum_x w(x)\gout_M(x)=\sum_\mu k_\mu$, and likewise for $\gin_M$; hence
$\tfrac12\sum_x w(x)(\gin_M+\gout_M)(x)=\sum_\mu k_\mu$.  Since
$\Delta_M=\gin_M-\gout_M$,
\[
  \sum_\mu k_\mu-\del(M)
  =\tfrac12\sum_x w(x)\Bigl[(d^++d^-)-|d^+-d^-|\Bigr](x)
  =\sum_x w(x)\min(d^+,d^-)(x),
\]
which is a sum of nonnegative terms, vanishing exactly under the stated
hypothesis.
\end{proof}

Equation~\eqref{eq:arrow-count} is the reason for calling $\del$ the electron
distance: it is the number of electrons the drawn arrows move.  It also
supplies the comparison promised in Remark~\ref{rem:vs-editing-distance}, since
the right-hand side of~\eqref{eq:arrow-count-defect} is nonnegative.

\begin{observation}
\label{obs:cancellation-parity}
The correction term in~\eqref{eq:arrow-count-defect} is even at every
pair-valued locus, where $w=2$.  Consequently
$\del(M)\equiv\sum_\mu k_\mu\pmod 2$ unless some \emph{radical} locus is
simultaneously a source and a target.
\end{observation}

\section{Parity}
\label{sec:parity}

Integrality now enters as a genuine hypothesis rather than a convenience.

\begin{definition}
\label{def:signature}
The \emph{radical signature} of an integral multiset $M$ is the pair
$(r^+,r^-)$ with
\begin{equation}
  r^+=\sum_{v}\max\bigl(\Delta_M(\rho_v),0\bigr),\qquad
  r^-=\sum_{v}\max\bigl(-\Delta_M(\rho_v),0\bigr) ,
  \label{eq:signature}
\end{equation}
the numbers of unpaired electrons created and destroyed.  We write
$\Delta R=r^+-r^-$ and $R_{\mathrm{abs}}=r^++r^-$.
\end{definition}

\begin{theorem}
\label{thm:parity}
Let $M$ be integral.  Then $\Delta R$ is even, and
\begin{equation}
  \del(M)\equiv r^+\equiv r^-\pmod 2 .
  \label{eq:parity}
\end{equation}
The electron distance is odd if and only if an odd number of unpaired
electrons is created, equivalently destroyed.
\end{theorem}

\begin{proof}
Write $L=\sum_v\Delta\ell_v$.  Expanding~\eqref{eq:conservation} in the four
sorts,
\begin{equation}
  2L+\Delta R+2\sum_e\Delta o_e=0 .
  \label{eq:conservation-expanded}
\end{equation}
By integrality $L$ and every $\Delta o_e$ lie in $\mathbb Z$, so reduction
of~\eqref{eq:conservation-expanded} modulo $2$ yields $\Delta R\equiv0$, that
is $r^+\equiv r^-$.  Solving~\eqref{eq:conservation-expanded} gives
$\sum_e\Delta o_e=-L-\Delta R/2$.  Reducing~\eqref{eq:distance-fields} modulo
$2$, using Observation~\ref{obs:sigma-pi} and $|a|\equiv a$ on the two
integer-valued sums,
\[
  \del\equiv\Bigl(-L-\tfrac{\Delta R}{2}\Bigr)+L+\tfrac{R_{\mathrm{abs}}}{2}
    \equiv\tfrac12\bigl(R_{\mathrm{abs}}-\Delta R\bigr)=r^-\pmod2 ,
\]
where $-x\equiv x\pmod 2$ was used.  With $r^+\equiv r^-$ this
is~\eqref{eq:parity}.
\end{proof}

\begin{remark}
\label{rem:heg-43}
The first conclusion of Theorem~\ref{thm:parity} recovers Eq.~(4.3) of~[1].
There, for a reaction whose atomwise discrepancies $\dsc_\alpha(v)$ all
vanish, one computes
$\sum_{v}\bigl(R'(\alpha(v))-R(v)\bigr)=2\bigl(|E|-|E'|\bigr)$ by observing
that every bond is incident with exactly two vertices, and concludes that the
total unpaired discrepancy is even.  In the locus formulation the same
statement is the mod-$2$ reduction of~\eqref{eq:conservation-expanded}, where
the factor $2$ appears not from double incidence but from the quota
$w(\sigma)=w(\pi)=w(\lp)=2$.  Theorem~\ref{thm:parity} strengthens the
conclusion from the parity of $\Delta R$, which is the parity of the
\emph{net} radical change, to the parity of $r^+$ and $r^-$ separately, which
is what controls $\del$.
\end{remark}

\begin{remark}
\label{rem:parity-needs-integrality}
Integrality cannot be dropped.  The original annotation of record~279 of the
radical corpus (Table~\ref{tab:s4}) has $\Delta o_{14,15}=-\tfrac32$, hence
signature $(r^+,r^-)=(2,1)$ with $\Delta R=1$ odd, which
Theorem~\ref{thm:parity} forbids, and $\del=4$, contradicting~\eqref{eq:parity}
against $r^-$.  Five of the six non-integral multisets catalogued in
Section~\ref{sec:audit} violate~\eqref{eq:parity}.
\end{remark}

\section{The polar and radical cases}
\label{sec:dichotomy}

Theorem~\ref{thm:parity} answers the parity question in terms of $r^\pm$.  We
now convert it into the two statements it was intended to supply.

\subsection{Polar groups}

\begin{theorem}
\label{thm:polar}
Let $M$ consist of two-electron moves only, so $f=0$.  Then
\begin{enumerate}
\item[(i)] $M$ is integral, $\Delta_M\in\mathbb Z^{\Loc(G)}$;
\item[(ii)] $\del(M)=2|M|$, and in particular $\del(M)$ is even.
\end{enumerate}
Neither conclusion requires locality, availability or $\Pi$-validity.
\end{theorem}

\begin{proof}
(i) A move $x\xrightarrow{2}y$ changes $c$ by $-2/w(x)$ at $x$ and $+2/w(y)$
at $y$; since $w\in\{1,2\}$ divides~$2$ in either case, both quantities are
integers, and $\Delta_M$ is a finite sum of integers.

(ii) Each move contributes $k_\mu=2$, so $\sum_\mu k_\mu=2|M|$.  If no locus
is both source and target, Proposition~\ref{prop:arrow-count}
gives $\del(M)=2|M|$ directly.  Otherwise, every $d^\pm_M(x)$ is an integer
multiple of $2/w(x)$, so $w(x)\min(d^+,d^-)(x)$ is an even integer at every
locus; by~\eqref{eq:arrow-count-defect}, $\del(M)=2|M|$ less an even number,
and is even.
\end{proof}

Assertion (i) is where the two-electron hypothesis is spent, and it is the
reason the argument does not extend: a fishhook contributes $\pm\tfrac12$ at a
pair-valued locus, and $1/w(x)$ is not an integer when $w(x)=2$.

\begin{corollary}
\label{cor:polar-mechanism}
Let $(M_1,\dots,M_n)$ be a mechanism all of whose moves are two-electron.
Then each $\del(M_i)$ is even, the step sum $\sum_i\del(M_i)$ is even, and the
net distance $\del\bigl(\sum_i\Delta_{M_i}\bigr)$ is even.  A polar reaction is
even at every level of description: per arrow, per step, per mechanism, and
end to end.
\end{corollary}

\begin{proof}
Each step is even by Theorem~\ref{thm:polar}(ii), hence so is the sum.  The
accumulated delta is integral by Theorem~\ref{thm:polar}(i) and has
$r^+=r^-=0$, since no move touches a radical locus with an odd quantum;
Theorem~\ref{thm:parity} gives $\del\equiv0$.
\end{proof}

\subsection{Radical groups}

\begin{theorem}
\label{thm:dichotomy}
Let $M$ be integral.  Then $\del(M)$ is odd if and only if $r^+$ is odd.
Integrality admits precisely the signatures with $r^+\equiv r^-\pmod2$, and
these are classified by Table~\ref{tab:signature}.
\end{theorem}

\begin{proof}
Both assertions are contained in Theorem~\ref{thm:parity}.
\end{proof}

\begin{table}[htbp]
\centering
\small
\caption{Parity of the electron distance by radical signature, for integral
move multisets.  Signatures with $r^+\not\equiv r^-$ are excluded by
Theorem~\ref{thm:parity}.}
\label{tab:signature}
\begin{tabular}{@{}c l c l@{}}
\toprule
$(r^+,r^-)$ & mechanistic reading & $\del$ & occurrences \\
\midrule
$(0,0)$ & closed shell, polar & even & the nine polar classes \\
$(2,0)$ & initiation, homolysis & even & M1; record 3{,}970 \\
$(0,2)$ & termination, recombination & even & M2; records 2{,}516, 2{,}300 \\
$(1,1)$ & propagation & \textbf{odd} & M3--M7; conj.\ abstraction \\
$(2,2)$ & concerted double propagation & even & --- \\
$(3,1),(1,3)$ & higher odd transfer & odd & --- \\
\midrule
\multicolumn{4}{@{}l}{$(1,0),(0,1),(2,1),\dots$ excluded: $r^+\not\equiv r^-$
contradicts integrality} \\
\bottomrule
\end{tabular}
\end{table}

\begin{corollary}
\label{cor:chain-carrier}
An integral radical group has odd electron distance precisely when it
relocates an odd number of unpaired electrons.  A group that creates or
destroys unpaired electrons in pairs---initiation and termination---is even,
as is any closed-shell group.  In the taxonomy of the companion manuscript
this partitions the macros: M1 and M2 are even with $\del=2$; M3, M4, M5, M6
and M7 are odd with $\del=3$; conjugated hydrogen abstraction is odd with
$\del=5$.
\end{corollary}

\begin{example}
\label{ex:beta-scission}
The $\beta$-scission of the neopentoxyl radical is the group
\[
  \rho_o\xrightarrow{1}\pi_{oc_1},\qquad
  \sigma_{c_1c_2}\xrightarrow{1}\pi_{oc_1},\qquad
  \sigma_{c_1c_2}\xrightarrow{1}\rho_{c_2},
\]
with $\Delta(\rho_o,\pi_{oc_1},\sigma_{c_1c_2},\rho_{c_2})=(-1,1,-1,1)$ and
signature $(1,1)$.  Hence $\del=3$, odd, in agreement with
Corollary~\ref{cor:chain-carrier} and with $2c+f=3$.  The electron-blind
distance is $\del_{\mathrm b}=2$: one $\pi$ formed and one $\sigma$ broken.
The discrepancy is precisely the relocation of the unpaired electron from $o$
to $c_2$, which a bond-only account does not register, and oddness of $\del$ is
thus the formal signature of a chain-carrying step.
\end{example}

\subsection{Counting fishhooks}

\begin{corollary}
\label{cor:fishhook}
If no radical locus of $M$ is both a source and a target, then
\begin{equation}
  \del(M)\equiv f\pmod 2 ,
  \label{eq:fishhook-parity}
\end{equation}
and if moreover $M$ is integral, $f\equiv r^+\equiv r^-\pmod2$.  Curly arrows
do not affect the parity of $\del$.
\end{corollary}

\begin{proof}
By Proposition~\ref{prop:arrow-count}, $\sum_\mu k_\mu=2c+f\equiv f$, and by
Observation~\ref{obs:cancellation-parity} the correction term is even at every
pair-valued locus; the excluded hypothesis is the remaining case.  The second
assertion follows with Theorem~\ref{thm:parity}.
\end{proof}

\begin{example}
\label{ex:fishhook-fails}
The hypothesis of Corollary~\ref{cor:fishhook} is not vacuous.  The integral
multiset with electron delta
\[
  E(\pi_e)=-2,\quad E(\sigma_e)=+6,\quad E(\lp_a)=-2,\quad E(\lp_b)=-2,\quad
  E(\rho_w)=+1,\quad E(\rho_v)=-1
\]
has $f=2$ and $\del=7$, so that~\eqref{eq:fishhook-parity} fails; here $\rho_w$
both receives and emits an electron.  Its signature is $(1,1)$, and
Theorem~\ref{thm:parity} holds.  The radical signature is therefore the robust
invariant and the fishhook count the convenient one.  No group occurring in
the companion manuscript triggers the exception.
\end{example}

\section{Composition along a mechanism}
\label{sec:composition}

By Proposition~\ref{prop:metric} the triangle inequality gives
$\del\bigl(\sum_i\Delta_{M_i}\bigr)\le\sum_i\del(M_i)$ for any mechanism.  The
following identifies the parity of the slack.

\begin{proposition}
\label{prop:composition}
Let $(M_1,\dots,M_n)$ be a mechanism with net delta
$\Delta=\sum_i\Delta_{M_i}$, all $M_i$ integral, and put
\begin{equation}
  \kappa=\tfrac12\sum_{v\in V}
    \Bigl[\sum_{i}\bigl|\Delta_{M_i}(\rho_v)\bigr|-\bigl|\Delta(\rho_v)\bigr|\Bigr]
  \;\in\;\mathbb N_0 ,
  \label{eq:kappa}
\end{equation}
the number of unpaired electrons created at one step and destroyed at another.
Then
\begin{equation}
  \sum_{i=1}^{n}\del(M_i)\equiv\del(\Delta)+\kappa\pmod 2 .
  \label{eq:composition}
\end{equation}
Parity is additive, $\sum_i\del(M_i)\equiv\del(\Delta)$, precisely when
$\kappa$ is even, and in particular when no radical is both produced and
consumed within the mechanism.
\end{proposition}

\begin{proof}
For integers $a_i$ one has $\sum_i|a_i|\equiv\sum_i a_i\equiv|\sum_i a_i|
\pmod2$, so each bracket in~\eqref{eq:kappa} is an even integer and $\kappa\in
\mathbb N_0$.  Apply the same congruence locuswise to
$\del=\tfrac12\sum_x w(x)|E(x)|$.  At a pair-valued locus $w=2$ and the
discrepancy $\tfrac12 w(x)\bigl[\sum_i|\Delta_i(x)|-|\Delta(x)|\bigr]$ equals
$\sum_i|\Delta_i(x)|-|\Delta(x)|$, which is even and vanishes modulo $2$.  At a
radical locus $w=1$ and the discrepancy is
$\tfrac12\bigl[\sum_i|\Delta_i(\rho_v)|-|\Delta(\rho_v)|\bigr]$; summing these
gives $\kappa$.
\end{proof}

\begin{example}
\label{ex:chain}
Consider initiation, hydrogen-atom transfer and termination,
\[
  \ce{A-A -> 2A^{.}},\qquad
  \ce{A^{.} + H-B -> A-H + B^{.}},\qquad
  \ce{B^{.} + A^{.} -> A-B},
\]
with signatures $(2,0)$, $(1,1)$ and $(0,2)$ and step distances $2$, $3$
and~$2$.  The step sum is $7$, odd.  The net delta has $\Delta r_v=0$ for every
$v$ and four unit bond changes, so $\del(\Delta)=4$, even.  Each of
$\rho_{A_1}$, $\rho_{A_2}$ and $\rho_B$ is produced once and destroyed once,
whence $\kappa=3$ and $7\equiv4+3\pmod2$ as
Proposition~\ref{prop:composition} requires.
\end{example}

Example~\ref{ex:chain} shows that for a radical mechanism the two distances
must be distinguished.  The step sum counts every electron the arrows move,
including the round trips of the chain carriers; the net distance sees only
the overall transformation, which is closed-shell and therefore even by
Theorem~\ref{thm:parity}.  For polar mechanisms
Corollary~\ref{cor:polar-mechanism} makes the distinction immaterial.

\section{Diagnostic value and its limits}
\label{sec:diagnostics}

Theorem~\ref{thm:parity} and Observation~\ref{obs:sigma-pi} supply two
arithmetic tests for non-integrality.  Both are one-directional, and neither
replaces the integrality gate.

\begin{corollary}
\label{cor:diag-fractional}
If $\del_{\mathrm b}(M)\notin\mathbb Z$ then $M$ is not integral.
\end{corollary}

\begin{corollary}
\label{cor:diag-parity}
If $\del(M)\not\equiv r^+\pmod2$, or if $r^+\not\equiv r^-\pmod 2$, then $M$
is not integral.
\end{corollary}

\begin{remark}
\label{rem:diag-limits}
Neither converse holds.  The original annotation of record~404
(Table~\ref{tab:s4}) has $\Delta o_{5,21}=\Delta o_{5,22}=+\tfrac12$, two
half-units carried by \emph{distinct} edges, so $\del_{\mathrm b}=2\in\mathbb Z$
and the signature is $(1,1)$ with $\del=3$ odd; the record satisfies both
tests and is not integral.
\end{remark}

More fundamentally, $\del$ is a function of $\Delta_M$ alone, and is therefore
blind to every admissibility condition that is not.

\begin{observation}
\label{obs:blind}
The electron distance cannot detect a failure of common-prestate availability,
since two multisets with equal $\Delta_M$ have equal $\del$ by definition;
nor a failure of locality, since retargeting a move to a local partner leaves
the multiset of locus sorts and hence, generically, $\del$ unchanged; nor
disagreement with a declared endpoint, by Example~\ref{ex:2207} below.
\end{observation}

\begin{example}
\label{ex:2207}
Record 2{,}207 of the radical corpus, the single case left unresolved in the
companion manuscript, has annotation
$9\rightharpoonup(9,13)$, $(12,13)\rightharpoonup(9,13)$,
$(12,13)\rightharpoonup12$, with
\[
  \Delta o_{9,13}=+1,\quad\Delta o_{12,13}=-1,\quad
  \Delta r_9=-1,\quad\Delta r_{12}=+1,\qquad \del=3 ,
\]
whereas its declared endpoint breaks $(14,15)$, forms $(9,15)$, consumes the
radical at map~$9$ and declares the product radical at map~$12$:
\[
  \Delta o_{14,15}=-1,\quad\Delta o_{9,15}=+1,\quad
  \Delta r_9=-1,\quad\Delta r_{12}=+1,\qquad \del=3 .
\]
The two agree in $\del$, in $\del_{\mathrm b}$, in signature and in parity, and
disagree at every bond locus.  Equality of electron distance is thus necessary
for an annotation to realize an endpoint, and not sufficient.
\end{example}

\section{Audit of the companion manuscript}
\label{sec:audit}

The tables below evaluate Definition~\ref{def:distance} on every case
occurring in the companion manuscript.  The column $\sum k$ reports the
pushed-electron count of Proposition~\ref{prop:arrow-count}, written $c$ for
curly arrows and $f$ for fishhooks; it agrees with $\del$ throughout, no group
having a locus that is both source and target.  An asterisk marks a
non-integral multiset.

\begin{table}[htbp]
\centering
\footnotesize
\setlength{\tabcolsep}{4pt}
\caption{Molecular examples.}
\label{tab:molecular}
\begin{tabular}{@{}l l c c c c@{}}
\toprule
Case & nonzero $\Delta_M$ & $\del_{\mathrm b}$ & $\del$ & $r^+$ & $\sum k$ \\
\midrule
\ce{NH3 + H+ -> NH4+}
  & $\ell_N{-}1,\ s_{NH}{+}1$ & $1$ & $2$ & $0$ & $1c=2$ \\
\ce{H2N- + H+ -> NH3}
  & $\ell_N{-}1,\ s_{NH}{+}1$ & $1$ & $2$ & $0$ & $1c=2$ \\
\ce{NH3 + HCl}
  & $\ell_N{-}1,\ s_{NH}{+}1,\ s_{HCl}{-}1,\ \ell_{Cl}{+}1$ & $2$ & $4$ & $0$ & $2c=4$ \\
$S_N2$, \ce{NH3 + CH3Cl}
  & $\ell_N{-}1,\ s_{NC}{+}1,\ s_{CCl}{-}1,\ \ell_{Cl}{+}1$ & $2$ & $4$ & $0$ & $2c=4$ \\
$\beta$-scission
  & $r_o{-}1,\ p_{oc_1}{+}1,\ s_{c_1c_2}{-}1,\ r_{c_2}{+}1$ & $2$ & $3$ & $1$ & $3f=3$ \\
peroxide homolysis
  & $s_{OO}{-}1,\ r_{O_1}{+}1,\ r_{O_2}{+}1$ & $1$ & $2$ & $2$ & $2f=2$ \\
formaldehyde $\pi\!\to\!\lp$
  & $p_{CO}{-}1,\ \ell_O{+}1$ & $1$ & $2$ & $0$ & $1c=2$ \\
\bottomrule
\end{tabular}
\end{table}

\begin{table}[htbp]
\centering
\footnotesize
\setlength{\tabcolsep}{5pt}
\caption{The nine polar classes.  Uniformity of $\del$ is
Theorem~\ref{thm:polar}(ii) with $|M|=1$; the classes are distinguished by
where the pair goes, not by how far.}
\label{tab:polar}
\begin{tabular}{@{}l l c c@{}}
\toprule
class & typical context & $\del_{\mathrm b}$ & $\del$ \\
\midrule
$\lp\to\lp$       & adjacent nonbonding-pair relocation & $0$ & $2$ \\
$\lp\to\pi$       & nucleophilic addition               & $1$ & $2$ \\
$\lp\to\sigma$    & $S_N2$, proton transfer             & $1$ & $2$ \\
$\pi\to\lp$       & ionization, retro-addition          & $1$ & $2$ \\
$\pi\to\pi$       & conjugated or pericyclic shift      & $2$ & $2$ \\
$\pi\to\sigma$    & electrophilic addition              & $2$ & $2$ \\
$\sigma\to\lp$    & heterolysis, E1                     & $1$ & $2$ \\
$\sigma\to\pi$    & elimination, retro-cycloaddition    & $2$ & $2$ \\
$\sigma\to\sigma$ & 1,2-migration, hydride transfer     & $2$ & $2$ \\
\bottomrule
\end{tabular}
\end{table}

\begin{table}[htbp]
\centering
\footnotesize
\setlength{\tabcolsep}{4pt}
\caption{The seven radical macros and the conjugated variant, illustrating
Corollary~\ref{cor:chain-carrier}.}
\label{tab:macros}
\begin{tabular}{@{}l l c c c c@{}}
\toprule
macro & nonzero $\Delta_M$ & $\del_{\mathrm b}$ & $\del$ & $(r^+,r^-)$ & parity \\
\midrule
M1 homolysis
  & $s{-}1,\ r_u{+}1,\ r_v{+}1$ & $1$ & $2$ & $(2,0)$ & even \\
M2 recombination
  & $r_u{-}1,\ r_v{-}1,\ s{+}1$ & $1$ & $2$ & $(0,2)$ & even \\
M3 radical addition
  & $r_R{-}1,\ s_{RC_1}{+}1,\ p_{C_1C_2}{-}1,\ r_{C_2}{+}1$ & $2$ & $3$ & $(1,1)$ & odd \\
M4 $\beta$-scission
  & $r_O{-}1,\ p_{OC_1}{+}1,\ s_{C_1C_2}{-}1,\ r_{C_2}{+}1$ & $2$ & $3$ & $(1,1)$ & odd \\
M5 radical resonance
  & $r_{C_1}{-}1,\ p_{C_1C_2}{+}1,\ p_{C_2C_3}{-}1,\ r_{C_3}{+}1$ & $2$ & $3$ & $(1,1)$ & odd \\
M6 H abstraction
  & $r_R{-}1,\ s_{RH}{+}1,\ s_{HC}{-}1,\ r_C{+}1$ & $2$ & $3$ & $(1,1)$ & odd \\
M7 alpha resonance
  & $\ell_d{-}1,\ r_d{+}1,\ r_a{-}1,\ \ell_a{+}1$ & $0$ & $3$ & $(1,1)$ & odd \\
\addlinespace[2pt]
conj.\ H abstraction
  & $r_R{-}1,\ s_{RH}{+}1,\ s_{HC_1}{-}1,$ & $4$ & $5$ & $(1,1)$ & odd \\
  & $p_{C_1C_2}{+}1,\ p_{C_2C_3}{-}1,\ r_{C_3}{+}1$ & & & & \\
\bottomrule
\end{tabular}
\end{table}

Macro M7 is the sharpest illustration of the refinement.  It edits no bond, so
$\del_{\mathrm b}=0$ and the step is invisible to a classical account, while
$\del=3$ records the three electrons relocated between a lone pair and a
radical.  The value is computed from the declared expansion, not from the
surface arrow; read literally as a single one-electron $\lp\to\lp$ move, the
surface notation would give $\del=1$ with half-integral delta.

\begin{remark}
\label{rem:conj-footnote}
The conjugated abstraction row is computed from its endpoint delta, which is
unambiguous, rather than from an arrow list.  The footnote to the macro table
of the companion manuscript states that the simple group ``expands with
$\sigma\to\pi$ and $\boldsymbol\cdot\to\pi$'', but the group realizing
$\ce{R^{.} + H-C_1-C_2=C_3 -> R-H + C_1=C_2-C_3^{.}}$ is
$\rho_R\to\sigma_{RH}$, $\sigma_{HC_1}\to\sigma_{RH}$,
$\sigma_{HC_1}\to\pi_{C_1C_2}$, $\pi_{C_2C_3}\to\pi_{C_1C_2}$,
$\pi_{C_2C_3}\to\rho_{C_3}$, whose added patterns relative to M6 are
$\sigma\to\pi$, $\pi\to\pi$ and $\pi\to\boldsymbol\cdot$, replacing
$\sigma\to\boldsymbol\cdot$.  The enumeration in that footnote appears to be
short one pattern and to name $\boldsymbol\cdot\to\pi$, which does not occur.
By Proposition~\ref{prop:metric} the value $\del=5$ is unaffected, since it
depends only on $\Delta_M$; the wording should nonetheless be checked.
\end{remark}

\footnotesize
\setlength{\tabcolsep}{4pt}
\begin{longtable}{@{}l >{\raggedright\arraybackslash}p{6.3cm} c c c c@{}}
\caption{Reviewed radical annotations, original and revised.  Bond loci are
written $(a,b)$; the source table does not resolve $\sigma$ from $\pi$, which
does not affect $\del$ since only $|\Delta o_e|$ enters.  An asterisk marks a
non-integral multiset.}
\label{tab:s4}\\
\toprule
record & nonzero $\Delta_M$ & $\del_{\mathrm b}$ & $\del$ & $(r^+,r^-)$ & $\sum k$ \\
\midrule
\endfirsthead
\toprule
record & nonzero $\Delta_M$ & $\del_{\mathrm b}$ & $\del$ & $(r^+,r^-)$ & $\sum k$ \\
\midrule
\endhead
\midrule
\multicolumn{6}{r@{}}{\small continued on next page}
\endfoot
\bottomrule
\endlastfoot
279 orig.$^\ast$
  & $r_2{-}1,\ o_{2,14}{+}1,\ o_{14,15}{-}\tfrac32,\ r_{15}{+}2$
  & $\tfrac52$ & $4$ & $(2,1)$ & $4f=4$ \\
279 rev.
  & $r_2{-}1,\ o_{2,14}{+}1,\ o_{14,15}{-}1,\ r_{15}{+}1$
  & $2$ & $3$ & $(1,1)$ & $3f=3$ \\
\addlinespace[2pt]
404 orig.$^\ast$
  & $r_5{-}1,\ o_{5,21}{+}\tfrac12,\ o_{21,22}{-}1,\ o_{5,22}{+}\tfrac12,\ r_{21}{+}1$
  & $2$ & $3$ & $(1,1)$ & $3f=3$ \\
404 rev.
  & $r_5{-}1,\ o_{5,22}{+}1,\ o_{21,22}{-}1,\ r_{21}{+}1$
  & $2$ & $3$ & $(1,1)$ & $3f=3$ \\
\addlinespace[2pt]
540 orig.$^\ast$
  & $r_1{-}1,\ o_{1,2}{+}\tfrac12$ & $\tfrac12$ & $1$ & $(0,1)$ & $1f=1$ \\
540 rev.
  & $\ell_1{-}1,\ o_{1,2}{+}1$ & $1$ & $2$ & $(0,0)$ & $1c=2$ \\
\addlinespace[2pt]
1{,}310
  & $r_{13}{-}1,\ o_{6,13}{+}1,\ o_{5,6}{-}1,\ r_5{+}1$ & $2$ & $3$ & $(1,1)$ & $3f=3$ \\
1{,}317
  & $r_{19}{-}1,\ o_{1,19}{+}1,\ o_{1,2}{-}1,\ r_2{+}1$ & $2$ & $3$ & $(1,1)$ & $3f=3$ \\
\addlinespace[2pt]
2{,}046 orig.
  & $r_{21}{-}1,\ r_{10}{+}1$ & $0$ & $1$ & $(1,1)$ & $1f=1$ \\
2{,}046 rev.
  & $r_{10}{-}1,\ r_{21}{+}1$ & $0$ & $1$ & $(1,1)$ & $1f=1$ \\
2{,}046 rev., M7 reading
  & $\ell_{10}{-}1,\ r_{10}{+}1,\ r_{21}{-}1,\ \ell_{21}{+}1$ & $0$ & $3$ & $(1,1)$ & $3f=3$ \\
\addlinespace[2pt]
2{,}207 annotation
  & $r_9{-}1,\ o_{9,13}{+}1,\ o_{12,13}{-}1,\ r_{12}{+}1$ & $2$ & $3$ & $(1,1)$ & $3f=3$ \\
2{,}207 endpoint
  & $o_{14,15}{-}1,\ o_{9,15}{+}1,\ r_9{-}1,\ r_{12}{+}1$ & $2$ & $3$ & $(1,1)$ & --- \\
\addlinespace[2pt]
2{,}300 orig.
  & $r_1{-}1,\ o_{1,2}{+}1,\ o_{2,3}{-}1,\ o_{3,4}{+}1,\ r_4{-}1$ & $3$ & $4$ & $(0,2)$ & $4f=4$ \\
2{,}300 rev.
  & $r_1{-}1,\ r_5{-}1,\ o_{1,5}{+}2,\ o_{1,2}{-}1,\ o_{2,4}{+}1,\ o_{4,5}{-}1$
  & $5$ & $6$ & $(0,2)$ & $2f{+}2c=6$ \\
\addlinespace[2pt]
2{,}516 orig.
  & $r_4{-}1,\ r_9{-}1,\ o_{4,8}{+}1$ & $1$ & $2$ & $(0,2)$ & $2f=2$ \\
2{,}516 rev.
  & $r_4{-}1,\ r_8{-}1,\ o_{4,8}{+}1$ & $1$ & $2$ & $(0,2)$ & $2f=2$ \\
\addlinespace[2pt]
3{,}970 orig.
  & $o_{20,21}{-}1,\ r_{20}{+}1,\ r_2{+}1$ & $1$ & $2$ & $(2,0)$ & $2f=2$ \\
3{,}970 rev.
  & $o_{20,21}{-}1,\ r_{20}{+}1,\ r_{21}{+}1$ & $1$ & $2$ & $(2,0)$ & $2f=2$ \\
\addlinespace[2pt]
4{,}852 orig.
  & $r_{10}{-}1,\ o_{10,20}{+}1,\ o_{20,21}{-}1,\ r_2{+}1$ & $2$ & $3$ & $(1,1)$ & $3f=3$ \\
4{,}852 rev.
  & $r_{10}{-}1,\ o_{10,20}{+}1,\ o_{20,21}{-}1,\ r_{21}{+}1$ & $2$ & $3$ & $(1,1)$ & $3f=3$ \\
\end{longtable}
\normalsize
\setlength{\tabcolsep}{6pt}

Three features of Table~\ref{tab:s4} deserve comment.  Every revision that
merely \emph{retargets} an arrow---records 404, 2{,}046, 2{,}516, 3{,}970 and
4{,}852---leaves $\del$ unchanged, confirming
Observation~\ref{obs:blind}: locality failures and missing atom maps are
invisible to the electron distance.  The two revisions that alter the electron
count of an arrow do change it, namely record 540, where a fishhook is replaced
by a paired move and $\del$ rises from $1$ to $2$, and record 279, where a
duplicated fishhook is removed and $\del$ falls from $4$ to $3$.  Finally,
record 2{,}046 is the one entry the source table leaves ambiguous: read as a
bare one-electron atom-to-atom arrow it gives $\del=1$, and read as the alpha
resonance shorthand of macro M7 it gives $\del=3$.  Both are odd, consistent
with Theorem~\ref{thm:parity}, but the committed deltas differ and the record
should be disambiguated.

\begin{table}[htbp]
\centering
\footnotesize
\setlength{\tabcolsep}{5pt}
\caption{Rejected controls.  All have integral $\del$ while being non-integral
multisets, illustrating Remark~\ref{rem:integrality-weaker}.}
\label{tab:controls}
\begin{tabular}{@{}l l c c l@{}}
\toprule
control & nonzero $\Delta_M$ & $\del_{\mathrm b}$ & $\del$ & rejected by \\
\midrule
$\lp\to\sigma$ made one-electron
  & $\ell_N{-}\tfrac12,\ s_{NH}{+}\tfrac12$ & $\tfrac12$ & $1$ & integrality \\
homolysis, one fishhook removed
  & $s_{OO}{-}\tfrac12,\ r_O{+}1$ & $\tfrac12$ & $1$ & integrality \\
single fishhook
  & $s_{uv}{-}\tfrac12,\ r_u{+}1$ & $\tfrac12$ & $1$ & integrality \\
equal net deltas, unequal prestate
  & identical & equal & equal & availability \\
\bottomrule
\end{tabular}
\end{table}

The last row of Table~\ref{tab:controls} is the essential one.  The
corresponding figure of the companion manuscript is constructed precisely so
that two groups share a net delta and differ in admissibility; they therefore
share $\del$ exactly, and no function of $\Delta_M$ can separate them.

\begin{remark}[Computational verification]
\label{rem:verification}
All entries of Tables~\ref{tab:molecular}--\ref{tab:controls} were recomputed
from~\eqref{eq:distance} by the script \code{cd\_check.py} accompanying this
note, which tested~\eqref{eq:conservation}, Theorem~\ref{thm:integer},
Proposition~\ref{prop:arrow-count} and Theorem~\ref{thm:parity} on each case.
Equation~\eqref{eq:conservation} and Theorem~\ref{thm:integer} hold in all $49$
cases, including the six non-integral ones, as their hypothesis-free
statements require; \eqref{eq:arrow-count} holds in all $48$ cases possessing
an arrow list; and~\eqref{eq:parity} holds in all $43$ integral cases and fails
in five of the six non-integral ones, the exception being record~404 discussed
in Remark~\ref{rem:diag-limits}.  Theorems~\ref{thm:polar}
and~\ref{thm:dichotomy} and Proposition~\ref{prop:composition} were in
addition tested on $20\,000$ to $40\,000$ pseudorandom multisets each.
\end{remark}

\section{Concluding remarks}
\label{sec:conclusion}

The electron distance is the weighted $\ell^1$ metric on locus occupancies,
and its arithmetic is completely determined.  It is integer-valued without
hypothesis, which makes it useless as a test for integrality of a move
multiset and useful as a summary of one.  It counts electrons rather than
operations, which separates it from the editing distance of~[1] and makes it a
closed-form quantity rather than a minimum over realizing sequences.  Its
parity is controlled by the radical signature alone, from which the polar and
radical cases follow: two-electron chemistry is even unconditionally and at
every level of description, while radical chemistry is even for initiation,
termination and closed-shell steps and odd for propagation.

Three limitations should be recorded.  First, all parity statements presuppose
integrality; below that threshold only Lemma~\ref{lem:conservation} and
Theorem~\ref{thm:integer} survive, and Remark~\ref{rem:parity-needs-integrality}
exhibits a corpus record where the failure is real rather than hypothetical.
Second, $\del$ is a function of the net delta, hence blind to common-prestate
availability, to locality, and to agreement with a declared endpoint; the
conformance layer of the companion manuscript is not replaceable by an
arithmetic invariant, and Example~\ref{ex:2207} exhibits an annotation and an
endpoint that agree in $\del$ and disagree everywhere else.  Third, the results
are stated for the two-centre Lewis model with the quota system
of~\eqref{eq:quota}; multicentre bonding or fractional occupation would
require different weights, and with them a different parity theory.

A question left open is whether the parity classification of
Table~\ref{tab:signature} admits a converse: whether every signature permitted
by Theorem~\ref{thm:parity} is realized by an admissible event group over
$\Pi_{\mathrm{exec}}$.  The signatures $(2,2)$, $(3,1)$ and $(1,3)$ are not
attested in either corpus examined here.

\begin{thebibliography}{9}
\bibitem{heg} R.~Holzschuh, C.~Flamm, D.~Merkle and P.~F.~Stadler,
  \emph{Half-edge graphs, Lewis formulas, and electron pushing}.
  ChemRxiv preprint, 13 July 2026.
  \href{https://doi.org/10.26434/chemrxiv.15006001/v1}{doi:10.26434/chemrxiv.15006001/v1}.
\bibitem{aam} M.~E.~Gonz\'alez Laffitte, T.-L.~Phan and P.~F.~Stadler,
  \emph{Extension of partial atom-to-atom maps: uniqueness and algorithms}.
  Algorithms for Molecular Biology \textbf{21} (2026).
\end{thebibliography}

\end{document}
